% Solutions to the exercises of Chapter 22.
\begin{solutions}
  \item For a straight edge with unit normal $n$, every gradient in the window is parallel to $n$, $\nabla I = g(x)\,n$, so
        \[
          M = \Big(\sum_xw(x)g(x)^2\Big)\,nn\T,
        \]
        of rank one: eigenvalues $\lambda_1 = \sum_xwg^2$ and $\lambda_2 = 0$, with the null eigenvector along the edge. For a vertical step of height $h$ in a $15\times15$ uniform window with
        central differences, each row has two pixels with $I_u = h/2$, so $M = \diag(7.5h^2, 0)$. For an ideal corner (a bright quadrant) with the window centered on
        its tip, each edge covers half the window: $M = h^2\big[\begin{smallmatrix} 4 & 0.25 \\ 0.25 & 4\end{smallmatrix}\big]$ with eigenvalues $3.75h^2$ and $4.25h^2$, both large; the small
        off-diagonal entry comes from the tip pixel, where both gradients are non-zero. More generally, two edges of equal strength $c$ with normals at an
        angle $\alpha$ give $M = c(n_1n_1\T + n_2n_2\T)$ with eigenvalues $c(1 \pm \cos\alpha)$: a right-angle corner is ideal, while a sharp wedge (normals $20^\circ$ apart,
        $\lambda_2 = 0.06c$) is poorly localized along its bisector.

  \item The characteristic polynomial of $M = \big[\begin{smallmatrix} a & b \\ b & c\end{smallmatrix}\big]$ is $\lambda^2 - (\tr M)\lambda + \det M$, so $\lambda = \tfrac12\big(\tr M \pm \sqrt{(\tr M)^2 - 4\det M}\big)$.
        The discriminant is
        \[
          (a + c)^2 - 4(ac - b^2) = (a - c)^2 + 4b^2 \ge 0,
        \]
        so both roots are real (as they must be for a symmetric matrix) and
        $\lambda_{\min} = \tfrac12\big(a + c - \sqrt{(a - c)^2 + 4b^2}\big)$. Use the second form in code: $(\tr M)^2 - 4\det M$ subtracts two nearly equal numbers when the eigenvalues
        are close, and rounding can make it slightly negative, which turns the square root into a NaN.

  \item On a texture smoothed with $\sigma = 2$~px (as in \cref{fig:lk}) and a $21\times21$ window, the shift $(0.3, -0.7)$ gives errors of $0.76$, $0.022$ and $0.010$~px after
        zero, one and two iterations, and then stalls at $0.010$~px. The floor is the bias of bilinear interpolation, since the shifted image was generated with
        cubic interpolation. With cubic sampling the error keeps falling, as in the figure. For a shift of 8~px and a $15\times15$ window, single-level tracking is correct to
        $0.1$~px from only 30--45\% of window positions: elsewhere it converges to a wrong local minimum. The texture decorrelates over a few pixels (the
        autocorrelation is $0.58$ at 3~px and $0.26$ at 5~px), so the linearization holds only for displacements of about 2--3~px. With one pyramid level above the
        base, the displacement there is 4~px and 93--98\% succeed; with two levels it is 2~px and all succeed. The rule: choose $L$ so that the largest expected
        displacement divided by $2^L$ is within the texture's correlation length.

  \item Let the true motion be $d$. The forward track returns $\hat d_f = d + e_f$ and the backward track, started at $x + \hat d_f$, returns $\hat d_b = -d + e_b$, where $e_f$ and $e_b$
        are the estimation errors of the two tracks. The round-trip error is $\hat d_f + \hat d_b = e_f + e_b$: the true motion cancels exactly, whatever its size. For a track
        that converged to the right minimum, only the causes of estimation error remain: image noise, interpolation and the convergence tolerance. Without them,
        both tracks would be exact. Moreover, $e_b \approx -e_f$: the backward window is centered on $x + \hat d_f$, so both directions compare the same pixel pairs of
        the same two images. In a simulation with noise of 1\% of the intensity range ($15\times15$ window), the forward tracking error has a median of $0.03$~px,
        but the round-trip error only $0.001$~px (99th percentile $0.005$~px), and $10^{-4}$~px without noise. The check tests self-consistency, not accuracy. Failures are a different matter: when half of the window is occluded in the
        second image, the median round-trip error is $1.7$~px, and 97\% exceed $0.1$~px. Real images add deformation, blur and brightness change, which
        raise the errors of correct tracks, so a threshold of a few tenths of a pixel keeps a wide margin while catching most failures. Better still, set it
        from a histogram of round-trip errors on a real sequence: correct tracks form a sharp peak near zero and failures a long tail.

  \item A rotation by $\theta = 0.05$~rad about the camera's $y$-axis turns the optical axis toward $+x_c$, so $R_{c_2c_1} = R_y(\theta)\T$. The feature at the principal point has
        bearing $(0, 0, 1)$, which becomes $R_y(\theta)\T(0, 0, 1)\T = (-\sin\theta, 0, \cos\theta)$, so
        \[
          u_2 = c_u - f_u\tan\theta = c_u - 22.95, \qquad v_2 = c_v :
        \]
        the feature moves 23~px to the left, the same as $KR_{c_2c_1}K\inv\tilde u_1$ gives. The prediction does not depend on depth; the neglected translation, 1~m/s
        over 50~ms at 5~m depth for example, adds only $f_u\cdot0.05/5 \approx 4.6$~px, which LK recovers easily from the predicted start.

  \item Divide the image into cells (for example $8\times6$ cells of $94\times80$~px for EuRoC). Compute the corner response once, with non-maximum suppression and
        one global threshold. Then take candidates strongest first, skipping any in a cell that has reached its quota or within the minimum distance of an
        existing track or an already chosen corner. Count the existing features per cell before detecting, so that only depleted cells receive new ones. To
        evaluate, report two numbers: the fraction of occupied cells, and the spread of the feature bearings (the standard deviations along the principal
        axes of their covariance, after undistortion). The script \code{scripts/solutions/ch22\_bucketing.py} does this for 150 corners, 10~px apart, with and
        without a quota of 4 per cell. It uses the first frame of \code{data/MH\_01\_easy} when the dataset is present. On a synthetic $752\times480$ image with one strongly
        textured region and weak texture elsewhere (one fifth of the contrast):
        \begin{center}\small
        \begin{tabular}{@{}l c c c@{}}
        \toprule
        selection & occupied cells & most corners in a cell & bearing spread (rad) \\
        \midrule
        strongest 150 & 46\% & 18 & 0.26 / 0.21 \\
        bucketed, at most 4 per cell & 96\% & 4 & 0.55 / 0.32 \\
        \bottomrule
        \end{tabular}
        \end{center}
        Without bucketing, every corner lies on the textured region and half of the view is empty. With bucketing, the corners cover nearly every cell, and
        the bearings spread twice as wide along the long axis of the image, which is what conditions the stacked Jacobian. Bucketing redistributes corners
        but cannot create texture: with a global threshold, cells whose texture is below it stay empty. That is the right outcome, since a corner too weak
        to pass the threshold is also too weak to track.

  \item A pool of fixed capacity $F_{\max}$ (at least the per-cell quota times the number of cells, plus room for tracks that are lost but not yet processed):
        \begin{itemize}
          \item an array of $F_{\max}$ track slots, each with a 64-bit ID, a status (free, new, tracked, lost), a count, and a fixed array of \code{kMaxClones}
                observations $\{$clone ID or timestamp, normalized $(x, y)\}$;
          \item a stack of free slot indices, so allocation and release are $O(1)$, and parallel fixed-capacity arrays of the current pixel positions and slot
                indices, which are what the LK call consumes;
          \item IDs that increase monotonically and are never reused, so stale references are detectable.
        \end{itemize}
        The history cannot overflow: when the oldest clone of a track leaves the window, the track is consumed (AGENTS.md F5), and any continuation starts a new
        history. If the pool is full, the simplest correct policy is to detect no new features until slots free up; the alternative is to evict the weakest
        (shortest or lowest-score) tracks. Either way the event is counted and logged, and the capacity is sized so that it does not happen in normal
        operation. Discarding a track loses information but never corrupts the estimate, unlike reusing one.
\end{solutions}
